\section{Weighted supports and transport of comparisons}
\label{sec:budgets}

The exact box solution is expensive to obtain by direct iteration. The
recursive solver will instead preserve its order relative to comparisons
with restricted supports. This section specifies that restriction and
proves the two mechanisms that keep it useful: a decrease of weighted
support budgets, and transport of comparisons when bounds or centers
change. No computed answer is required to be a subsolution or a
supersolution itself.

Separate budgets for the two sides parallel the two dominion-size
precisions in Parys's recursion~\cite[Algorithm~2]{Parys2019} and the
recursive refinements presented by Lehtinen, Parys, Schewe, and
Wojtczak~\cite[Sections~3--4]{LehtinenParysScheweWojtczak2022}.
Here the constrained objects are weighted supports of full comparison
vectors, and the decrease rules are proved directly.

Choose positive rational weights
$\alpha,\beta\in\mathbb Q_{>0}^V$ and positive rational budgets $k,m$.
For a vertex set $S$, write
\[
 \alpha(S)=\sum_{i\in S}\alpha_i,\qquad
 \beta(S)=\sum_{i\in S}\beta_i.
\]
\begin{definition}[Admissible comparisons]\label{def:admissible-comparison}
A subsolution $y$ of $[a,b]$ is \emph{admissible} if
$\alpha(S_y)\le k$. A supersolution $z$ of that box is admissible if
$\beta(S_z)\le m$. The supports are relative to $[a,b]$ and have the
boundary conventions of Definition~\ref{def:box-comparisons}.
\end{definition}

Given these parameters and a reference width $1\le B\le B_*$, the
procedure $\Solve$ will always return an integer vector $p\in[a,b]$.
Its target comparisons are
\begin{equation}\label{eq:comparison-guarantees}
 \Prb(p\not\ge y)\le\frac1{16},\qquad
 \Prb(p\not\le z)\le\frac1{16},
\end{equation}
separately for each fixed admissible $y$ and each fixed admissible $z$.
The event $p\not\ge y$ means that the inequality fails in at least one
coordinate. Thus each success is a whole-vector event, but the comparison
is fixed outside the probability. Section~\ref{sec:solving} proves this
pointwise assertion; it does not assert simultaneous success for all
admissible comparisons.

\subsection{A rank that decreases with the support parameters}

Define the positive rational parameter $R$ and a probability measure
$\mu$ on $V$ by
\begin{equation}\label{eq:rank-parameter}
 R=km\sum_{i\in V}\frac1{\alpha_i\beta_i},\qquad
 \mu(S)=
 \frac{\displaystyle\sum_{i\in S}1/(\alpha_i\beta_i)}
      {\displaystyle\sum_{i\in V}1/(\alpha_i\beta_i)}.
\end{equation}
The proof will induct on the integer
\[
 d(R)=\min\{d\in\Z_{\ge0}:(9/10)^dR<1\}.
\]
The algorithm itself need only compare rational numbers. Whenever
$R\ge1$ and $R'\le(9/10)R$, we have
\[
 (9/10)^{d(R)-1}R'\le(9/10)^{d(R)}R<1,
\]
so the integer rank decreases by at least one.

\begin{lemma}[Small rank]\label{lem:small-rank}
If $R<1$, the vector
\[
 p_i=\begin{cases}b_i,&\beta_i>m,\\a_i,&\beta_i\le m\end{cases}
\]
satisfies $y\le p\le z$ for every admissible subsolution $y$ and every
admissible supersolution $z$ of $[a,b]$.
\end{lemma}

\begin{proof}
Each summand $km/(\alpha_i\beta_i)$ is smaller than one. Thus
$\alpha_i\beta_i>km$: the two conditions $\alpha_i\le k$ and
$\beta_i\le m$ cannot both hold. If $\beta_i>m$, an admissible
supersolution has no support at $i$ and must equal $b_i=p_i$, which is
also at least $y_i$. Otherwise $\alpha_i>k$, so every admissible
subsolution equals $a_i=p_i$ there, which is at most $z_i$.
This coordinatewise reasoning includes collapsed intervals and proves
the full assertion.
\end{proof}

Put $\eta=4/5$. Multiplying either budget by $\eta$ reduces $R$ by
that factor. We will also reduce rank when a support is concentrated on
a set having small $\mu$-measure.

\begin{lemma}[Reweighting]\label{lem:reweighting}
Let $H\subseteq V$ satisfy $\mu(H)\le1/2$, and keep the box $[a,b]$
fixed.
\begin{enumerate}[label=\textup{(\roman*)}]
\item Set $\alpha'_i=\alpha_i/2$ on $H$ and $\alpha'_i=\alpha_i$ off
$H$, and set $k'=3k/5$, leaving $\beta,m$ unchanged. Every admissible
subsolution $y$ with $\alpha(S_y\cap H)>\eta k$ remains admissible.
Every formerly admissible supersolution remains admissible as well.
\item Set $\beta'_i=\beta_i/2$ on $H$ and $\beta'_i=\beta_i$ off $H$,
and set $m'=3m/5$, leaving $\alpha,k$ unchanged. Every admissible
supersolution $z$ with $\beta(S_z\cap H)>\eta m$ remains admissible,
as does every formerly admissible subsolution.
\end{enumerate}
In either case,
\[
 R'=\frac35(1+\mu(H))R\le\frac9{10}R.
\]
\end{lemma}

\begin{proof}
For the first change, the new support weight is
\[
 \alpha'(S_y)=\alpha(S_y)-\tfrac12\alpha(S_y\cap H)
 <k-\tfrac12\eta k=\tfrac35k=k'.
\]
No supersolution parameter changes. Write
$A=\sum_i1/(\alpha_i\beta_i)$. Halving $\alpha$ on $H$ replaces this
sum by
\[
 A+\sum_{i\in H}\frac1{\alpha_i\beta_i}=A(1+\mu(H)).
\]
Multiplying by the new budget product $3km/5$ gives the formula for
$R'$. For the second change, the identical support calculation reads
\[
 \beta'(S_z)=\beta(S_z)-\tfrac12\beta(S_z\cap H)
 <m-\tfrac12\eta m=\tfrac35m=m'.
\]
The subsolution parameters are now unchanged, and the reciprocal-product
sum again gains the factor $1+\mu(H)$. Finally $\mu(H)\le1/2$ gives
the rank contraction, including equality at the endpoint.
\end{proof}

\subsection{Moving boundaries and recentering}

A change of box requires a new comparison vector, not just a new
interval. The following constructions verify its inequalities and give
its exact support. Support containment will preserve every relevant
budget even when a centered child box extends outside its parent.

\begin{lemma}[Boundary projections]\label{lem:boundary-projections}
Let $y$ be a subsolution and $z$ a supersolution of $[a,b]$.
\begin{enumerate}[label=\textup{(\roman*)}]
\item If $a\le a'\le b$, then $y'=\max(y,a')$ is a subsolution of
$[a',b]$ with $S_{y'}=\{i:y_i>a'_i\}\subseteq S_y$.
\item If $a\le b'\le b$, then $z'=\min(z,b')$ is a supersolution of
$[a,b']$ with $S_{z'}=\{i:z_i<b'_i\}\subseteq S_z$.
\end{enumerate}
These projections preserve admissibility for unchanged weights and
budgets. Lowering only the upper bound also preserves any subsolution
that still fits, with unchanged support; raising only the lower bound
preserves any supersolution that still fits, with unchanged support.
The new bounds need not satisfy an inequality involving $F$.
\end{lemma}

\begin{proof}
For the first projection, $a'\le y'\le b$. It is active exactly where
$y_i>a'_i$. At such a coordinate $y'_i=y_i>a'_i\ge a_i$, and hence
\[
 y'_i=y_i\le F_i(y)\le F_i(y').
\]
The first inequality is the old active comparison; the second uses
$y'\ge y$. This proves both the new comparison and support containment.
At an inactive coordinate no inequality is required.

For the second projection, $a\le z'\le b'$. Its active coordinates are
exactly those where $z_i<b'_i$, and at each of them $z'_i=z_i<b'_i\le b_i$.
Thus
\[
 z'_i=z_i\ge F_i(z)\ge F_i(z'),
\]
using the old supersolution inequality and $z'\le z$. Again the support
can only shrink. Positive weights then give admissibility on both sides.
Finally, only the lower bound defines a subsolution's support and only
the upper bound defines a supersolution's support. Altering the other
bound affects containment alone, which proves the remaining statements,
including degenerate new intervals.
\end{proof}

\begin{lemma}[Centered projections]\label{lem:center-projections}
Let $q\in[a,b]$ be an integer vector and $A$ a positive integer.
\begin{enumerate}[label=\textup{(\roman*)}]
\item For a subsolution $y$ of $[a,b]$, if
$d=\max_i(y_i-q_i)\ge2A$, then
\[
 y'=\max(y-d+A,q-A)
\]
is a subsolution of $[q-A,q+A]$ with support
\[
 S_{y'}=\{i:y_i-q_i>d-2A\}\subseteq S_y.
\]
\item For a supersolution $z$ of $[a,b]$, if
$e=\max_i(q_i-z_i)\ge2A$, then
\[
 z'=\min(z+e-A,q+A)
\]
is a supersolution of $[q-A,q+A]$ with support
\[
 S_{z'}=\{i:q_i-z_i>e-2A\}\subseteq S_z.
\]
\end{enumerate}
Each projected comparison is admissible when the parent's comparison
is admissible and the support weights and budget are unchanged.
\end{lemma}

\begin{proof}
The definition of $d$ gives $y-d+A\le q+A$, so $y'$ lies in the
centered box. Its active-coordinate test is
\[
 y'_i>q_i-A\quad\Longleftrightarrow\quad y_i-q_i>d-2A.
\]
Since $d-2A\ge0$, every active coordinate has $y_i>q_i\ge a_i$ and
belongs to the original support. At that coordinate, translation and
monotonicity yield
\[
 y'_i=y_i-d+A\le F_i(y)-d+A
       =F_i(y-d+A)\le F_i(y').
\]
Although $y-d+A$ need not be in the parent box, $F$ is defined on all
integer vectors, and the old comparison was used only at an original
support coordinate.

\begin{samepage}
Similarly $z+e-A\ge q-A$, so $z'$ is in the new box. The active test is
\[
 z'_i<q_i+A\quad\Longleftrightarrow\quad q_i-z_i>e-2A.
\]
\end{samepage}
Such a coordinate has $z_i<q_i\le b_i$ and is in $S_z$. There
\[
 z'_i=z_i+e-A\ge F_i(z)+e-A
       =F_i(z+e-A)\ge F_i(z').
\]
This proves both support formulas and comparison properties. Positivity
of the weights then transfers the budgets.
\end{proof}

There is a second consequence of support containment that will be used
in orientation. For every $H\subseteq V$,
\[
 \alpha(S_y\cap H)\ge\alpha(S_{y'}\cap H),\qquad
 \beta(S_z\cap H)\ge\beta(S_{z'}\cap H).
\]
Thus a strict support-weight lower bound proved for a projected
comparison also holds for its parent. With unchanged parameters,
$\mu(H)$ is unchanged. Neither conclusion requires one centered box
to lie inside the other or their boundary heights to be close.

\subsection{Detecting a threatened truncation}

Before truncating a box, the algorithm will test whether a comparison
near that boundary can be handled with a reduced budget. A translated
comparison gives the test: if successful, it forces enough motion to
continue tightening. In this lemma take any integer $s\ge1$, put $h=1/s$,
and let $B$ be the reference width. Their algorithmic values are fixed in
Section~\ref{sec:algorithm}.

\begin{lemma}[Tightening projection]\label{lem:tightening-projection}
Suppose $B>128s$ and define
\[
 \delta=\floor{B/(4s)},\qquad C=B-\delta,\qquad h=1/s.
\]
Then $\delta>0$, and the following statements hold for $[a,b]$.
\begin{enumerate}[label=\textup{(\roman*)}]
\item If $y$ is a subsolution with $d=\max_i(y_i-a_i)>C$, then
$u=\max(y-d+\delta,a)$ is a subsolution in the same box, $u\le y$, and
\begin{align*}
 S_u&=\{i:y_i-a_i>d-\delta\}
       \subseteq\{i:y_i\ge a_i+B-hB\},\\
 \max_i(u_i-a_i)&=\delta.
\end{align*}
\item If $z$ is a supersolution with $e=\max_i(b_i-z_i)>C$, then
$v=\min(z+e-\delta,b)$ is a supersolution in the same box, $v\ge z$, and
\begin{align*}
 S_v&=\{i:b_i-z_i>e-\delta\}
       \subseteq\{i:z_i\le b_i-B+hB\},\\
 \max_i(b_i-v_i)&=\delta.
\end{align*}
\end{enumerate}
\end{lemma}

\begin{proof}
We have $B/(4s)>32$, so $\delta\ge32$, whereas $\delta\le B/4$.
In particular $B-2\delta\ge B/2>0$ and $2\delta\le B/(2s)\le hB$.
For part (i), set $c=d-\delta$. The threat $d>B-\delta$ implies
\begin{equation}\label{bud:tightening-shift}
 c>B-2\delta\ge B-hB,\qquad c>0.
\end{equation}
It follows that $a\le u=\max(y-c,a)\le y\le b$. A coordinate is
active precisely when $y_i-a_i>c$. At such a coordinate $u_i=y_i-c$
and $y_i>a_i$, whence
\[
 u_i=y_i-c\le F_i(y)-c=F_i(y-c)\le F_i(u).
\]
Equation~\eqref{bud:tightening-shift} also places every active coordinate
strictly above $a_i+B-hB$, and therefore in the asserted inclusive
set; no rounding of that real threshold is used. Finally,
\[
 u_i-a_i=\max\{y_i-a_i-d+\delta,0\}\le\delta,
\]
with equality at a coordinate attaining $d$.

For part (ii), put $c=e-\delta$. The same inequalities show
$c>B-2\delta\ge B-hB$ and $c>0$, so $a\le z\le v=\min(z+c,b)\le b$.
Its active test is $b_i-z_i>c$. At an active coordinate $z_i<b_i$,
and
\[
 v_i=z_i+c\ge F_i(z)+c=F_i(z+c)\ge F_i(v).
\]
The active coordinate lies strictly below $b_i-B+hB$, giving the low
set in the statement. Moreover,
\[
 b_i-v_i=\max\{b_i-z_i-e+\delta,0\}\le\delta,
\]
with equality wherever $e$ is attained. This proves the second part.
\end{proof}

If the high set in part (i) has $\alpha$-weight at most $\eta k$, then
$u$ qualifies for that reduced budget. Any answer $r\ge u$ satisfies
$r_i-a_i\ge\delta$ somewhere and therefore triggers the lower-bound
update. The dual assertion applies to $v$ when the low set has
$\beta$-weight at most $\eta m$: an answer $r\le v$ forces an
upper-bound update by at least $\delta$. These witnesses depend only
on the current box and the protected comparison. Consequently they
can be fixed before the fresh random bits of the test call, as required
by the conditional proof in Section~\ref{sec:solving}.
